% Auto-generated from the current Wareeny thesis PDF/text sources.
% This file prioritizes structured text conversion into LaTeX.
\documentclass[12pt,a4paper]{report}
\usepackage[utf8]{inputenc}
\usepackage[T1]{fontenc}
\usepackage{geometry}
\usepackage{setspace}
\usepackage{hyperref}
\usepackage{enumitem}
\usepackage{graphicx}
\usepackage{titlesec}
\usepackage{parskip}
\usepackage{array}
\usepackage{tabularx}
\usepackage{xurl}
\usepackage{mathptmx}
\geometry{margin=1in}
\setstretch{1.2}
\renewcommand{\arraystretch}{1.2}
\setcounter{secnumdepth}{3}
\setcounter{tocdepth}{3}
\hypersetup{colorlinks=true,linkcolor=black,urlcolor=blue}
\urlstyle{same}
\newcommand{\chaptercover}[1]{%
  \clearpage
  \thispagestyle{empty}
  \vspace*{\fill}
  \begin{center}
    {\Huge\bfseries #1\par}
    \vspace{0.8cm}
    \rule{0.35\textwidth}{1pt}\par
    \vspace{0.6cm}
    {\Large Wareeny Thesis\par}
  \end{center}
  \vspace*{\fill}
  \clearpage
}
\begin{document}

\begin{titlepage}
\centering
{\Large Future University in Egypt\par}
\vspace{0.5cm}
{\large Faculty of Computers \& Information Technology\par}
{\large Department of Computer Science\par}
\vspace{2cm}
{\Huge\bfseries Wareeny\par}
\vspace{0.6cm}
{\Large An AI-Powered Mobile Assistive System for Visually Impaired Users\par}
\vspace{1.5cm}
{\large Thesis Submitted in Partial Fulfillment of the Requirements for the Degree of Bachelor of Science in Computer Science\par}
\vspace{1.5cm}
\begin{tabular}{c}
Youssef Amr -- 20221949 \\
Mohamed Moataz -- 20223394 \\
Omar Ahmed Said -- 20225281 \\
Ahmed Mohamed -- 20223031 \\
Mohamed Wael -- 20223960 \\
Mohamed Khaled -- 20224727 \\
\end{tabular}
\vfill
\begin{tabular}{l}
Professor: Prof. Abeer El Korany \\
Teaching Assistant: Linah Mohamed \\
\end{tabular}
\vfill
{\large 2025--2026\par}
\end{titlepage}

\pagenumbering{roman}
\chapter*{Abstract}
\addcontentsline{toc}{chapter}{Abstract}


This thesis presents Wareeny, an AI-assisted mobile application intended to support blind and low-vision users through camera-based perception, spoken feedback, and accessible interaction patterns. In its current implemented form, Wareeny is a Flutter-based mobile system organized around a mode-driven interface that includes reading, navigation, stuff-finder, and settings flows. The reading module combines still-image capture with a routed vision-language pipeline, while the navigation and object-finding modules rely on on-device YOLO-based streaming detection. The application also provides bilingual English-Arabic interaction, text-to-speech output, haptic feedback, and persisted user preferences for language, speech rate, and haptics. Architecturally, the system separates presentation, feature controllers, shared AI services, backend integration, and lightweight persistence, producing a codebase that is both assistive in purpose and extensible in structure. The implemented repository therefore represents a credible functional foundation for an accessibility-oriented mobile assistant, while broader safety workflows, richer Arabic-domain evaluation, and large-scale user studies remain future work.



\chapter*{List of Abbreviations}
\addcontentsline{toc}{chapter}{List of Abbreviations}
\begin{description}[style=nextline,leftmargin=2.5cm]
  \item[AI] Artificial Intelligence
  \item[API] Application Programming Interface
  \item[COCO] Common Objects in Context
  \item[CPU] Central Processing Unit
  \item[CV] Computer Vision
  \item[GAN] Generative Adversarial Network
  \item[GPU] Graphics Processing Unit
  \item[LF-VQA] Long-Form Visual Question Answering
  \item[LMM] Large Multimodal Models
  \item[MENA] Middle East and North Africa
  \item[OCR] Optical Character Recognition
  \item[RAM] Random Access Memory
  \item[TTS] Text-to-Speech
  \item[VQA] Visual Question Answering
  \item[WHO] World Health Organization
  \item[YOLO] You Only Look Once
\end{description}

\tableofcontents

\clearpage
\pagenumbering{arabic}


\chaptercover{Chapter 1\\Introduction}
\chapter{Introduction}

This chapter introduces the project, provides background context, defines the problem, and outlines the objectives and significance of the Wareeny assistive system for visually impaired users.

\section{Background of the Study}

This project focuses on assistive technology for blind and low-vision users using computer vision, Optical Character Recognition (OCR), audio, and haptic feedback. Smartphones now include powerful cameras, sensors, and on-device AI capabilities. These capabilities create an opportunity to deliver real-time environmental awareness without specialized hardware. Vision-impaired users face daily barriers both indoors and outdoors. Objects, obstacles, text, and signals are often inaccessible without assistance. Human assistance is not always available, and traditional aids do not provide contextual understanding. Current assistive applications handle isolated tasks. Many focus only on OCR or basic object detection. They lack context awareness, struggle in cluttered environments, and often overwhelm users with unnecessary information. There is a gap in integrated systems that combine object detection, text understanding, safety alerts, and directional feedback. Few systems filter information based on user intent or risk, and privacy handling remains limited. This study addresses the need for a unified mobile system that delivers accurate, timely, and safe assistance tailored to real-world situations. By leveraging modern deep learning architectures and mobile computing capabilities, the proposed system aims to significantly improve independence and safety for visually impaired individuals.

\section{Problem Statement}

Blind and low-vision users lack a single reliable mobile system that provides real-time contextual understanding of their surroundings. Existing solutions handle isolated functions and fail in complex environments. They do not prioritize safety-critical information and offer limited privacy control. This results in reduced independence and increased risk during daily activities.

\section{Aim and Objectives}

\subsection{General Aim}

The general aim of this project is to build a mobile assistive system that supports blind and low-vision users through real-time vision, audio, and haptic feedback. The system seeks to improve safety, independence, and access to information in both indoor and outdoor spaces.

\subsection{Specific Objectives}

The specific objectives of this project are:

\begin{enumerate}
  \item Detect indoor objects and estimate distance with high accuracy.
  \item Detect the object held in the user’s hand with low false detection rate.
  \item Read currency and identify value quickly.
  \item Detect outdoor obstacles and traffic lights with real-time response with very low latency.
  \item Read text and answer user questions from documents with correct intent recognition.
  \item Trigger danger mode and emergency call very quickly.
\end{enumerate}

\section{Scope of the Project}

\subsection{Included Features}

The following features are included in the scope of this project:

\begin{itemize}
  \item Indoor object detection, brand recognition, and in-hand detection.
  \item Outdoor object detection, obstacle awareness, and object finding.
  \item OCR for documents, signs, and labels.
  \item Integrating LMM-based question answering.
  \item Audio cues and haptic alerts.
  \item Emergency call and danger mode activation.
\end{itemize}

\subsection{Excluded Features}

The following aspects are excluded from the scope of this project:

\begin{itemize}
  \item Hardware Manufacturing: The project focuses on a software-based mobile application. It does not include the development or sale of custom wearable hardware (such as smart glasses).
\end{itemize}

\subsection{Constraints}

The project is subject to the following constraints:

\paragraph*{Data Constraints}

\begin{itemize}
  \item Limited availability of labeled datasets for in-hand detection, brand recognition, and currency.
  \item Dataset bias toward common objects and environments.
  \item Language and currency coverage restricted to selected regions.
\end{itemize}

\paragraph*{Hardware Constraints}

\begin{itemize}
  \item Mobile phone CPU, GPU, and memory limitations.
  \item Camera resolution and low-light performance vary across devices.
  \item Battery consumption restricts continuous real-time processing.
\end{itemize}

\paragraph*{Time Constraints}

\begin{itemize}
  \item Fixed academic timeline limits model training and optimization cycles.
  \item Restricted time for large-scale user testing.
  \item Limited duration for dataset collection and annotation.
\end{itemize}

\section{Significance of the Project}

\subsection{Beneficiaries}

The project benefits multiple stakeholders:

\begin{itemize}
  \item Blind and low-vision users: Gain independence in daily activities.
  \item Caregivers: Gain safety assurance for their dependents.
  \item Researchers: Gain access to a modular assistive vision system for further study.
  \item Organizations: Gain a deployable accessibility tool for their constituents.
\end{itemize}

\subsection{Potential Impact}

The potential impact of this project includes:

\begin{itemize}
  \item Reduces daily dependency on human assistance.
  \item Combines vision, language, audio, and haptics in one integrated system.
  \item Focuses on context filtering to reduce cognitive noise.
  \item Helps blind individuals perform daily tasks with greater ease and confidence.
\end{itemize}

\section{Project Methodology}

\subsection{Development Approach}

The project follows an iterative Agile development approach with the following characteristics:

\begin{itemize}
  \item Incremental feature delivery with regular milestones.
  \item Continuous testing after each module integration.
  \item User-centered design with feedback-driven refinement.
\end{itemize}

\subsection{Key Technologies and Tools}

The following technologies and tools are employed in the development:

\begin{itemize}
  \item Deep learning models for object detection and OCR.
  \item Mobile frameworks for on-device inference.
  \item Text recognition and assistant LMM models.
  \item Audio output and vibration APIs for feedback.
  \item Smartphone camera, microphone, and motion sensors.
\end{itemize}

\section{Project Organization}

This thesis is organized as follows:

\begin{itemize}
  \item Chapter One: Introduces the project, presents the problem statement, defines the aim and objectives, and outlines the scope.
  \item Chapter Two: Provides a literature review of assistive vision systems, object detection, OCR, and audio feedback technologies.
  \item Chapter Three: Presents the system analysis and design, including models, datasets, and design decisions.
\end{itemize}


\chaptercover{Chapter 2\\Literature Review}
\chapter{Literature Review}

\section{Introduction}

Visual impairment represents a significant global health challenge, affecting approximately 2.2 billion people worldwide, with over 1 billion cases that could have been prevented or remain unaddressed [19]. These individuals face substantial barriers in navigating their environments, accessing textual information, and maintaining independent daily activities. The burden of visual impairment falls disproportionately on developing regions, particularly in the Middle East and North Africa (MENA), where access to assistive technologies remains limited. Recent advances in artificial intelligence and edge computing have opened new possibilities for developing sophisticated assistive technologies. However, the market remains dominated by solutions designed with a Westerncentric focus, lacking robust support for non-Latin scripts, regional entities, and cultural contexts. In Egypt and the broader Arab world, visually impaired individuals encounter a critical scarcity of assistive technologies that are linguistically adapted for Arabic and culturally tailored to recognize region-specific currency, products, and environmental elements. Furthermore, the majority of existing commercial solutions rely heavily on cloud-based processing architectures, raising significant privacy concerns regarding the transmission of sensitive personal data, including images of users’ environments and biometric information. This project, Wareeny, proposes a novel, comprehensive mobile application designed to bridge these specific accessibility gaps. By integrating state-of-the-art 2025 AI architectures—combining real-time AI processing for safety with Large Multimodal Models (LMMs) for semantic understanding—Wareeny aims to empower visually impaired users in the Arab world with a reliable, localized, and privacy-preserving assistive tool.

\section{Review of Related Systems and Studies}

The domain of assistive technology for the visually impaired is characterized by a mix of commercial mobile applications and academic research targeting specific computer vision problems. A critical review of these systems highlights both the progress made and the significant limitations that persist.

\subsection{Analysis of Commercial Assistive Applications}

Several prominent mobile applications currently serve the visually impaired community. However, a detailed analysis reveals deficiencies in accuracy, localization, and usability, particularly for users in the MENA region.

\paragraph*{Microsoft Seeing AI (2017)}

Microsoft Seeing AI [7], launched in 2017, represents one of the earliest comprehensive AI-powered assistive applications, providing multiple channels for text reading, scene description, and object identification. The application leverages cloud-based computer vision models to deliver real-time audio feedback through smartphone cameras. However, systematic user studies have revealed several critical limitations. First, the system exhibits reduced accuracy in complex, cluttered environments, where it struggles to distinguish foreground objects from background elements. Second, users report experiencing cognitive overload due to the system’s tendency to narrate non-essential environmental details rather than prioritizing safety-critical information such as obstacles or hazards. Third, its natural language processing capabilities are optimized predominantly for English, with substantially degraded performance when processing Arabic text due to the morphological complexity and cursive nature of Arabic script. Finally, the application’s exclusive availability on iOS limits accessibility for Android users, particularly in developing regions where Android devices dominate the market.

\paragraph*{Google Lookout (2019)}

Google Lookout [8], released in 2019, employs camera-based computer vision to identify objects, text, and people in the user’s environment. The application offers three operational modes: Explore, Shopping, and Scan, each optimized for specific use cases. Despite these capabilities, the system faces multiple technical and accessibility limitations. First, it is exclusively available on the Android platform, effectively excluding the substantial iOS user base. Second, users have reported system instability characterized by frequent application crashes, particularly during extended usage sessions or when processing complex visual scenes. Third, the continuous camera streaming and real-time AI inference impose substantial computational demands, resulting in rapid battery depletion that limits practical usage duration. Most critically, the application relies predominantly on cloud-based processing, requiring transmission of captured images to Google’s servers for analysis. This architecture introduces privacy vulnerabilities, particularly when users scan sensitive personal documents, financial information, or private correspondence. Additionally, the cloud dependency introduces latency in response times and renders the application non-functional in offline scenarios or areas with limited network connectivity.

\paragraph*{Envision AI (2018)}

Envision AI [9], introduced in 2018, has gained recognition for its optical character recognition (OCR) capabilities and cross-platform availability on both iOS and Android devices. The application provides functionality for document reading, object recognition, scene description, and color identification. However, it exhibits significant limitations in regional localization and technical architecture. The system’s object detection models are trained primarily on Western-centric datasets such as COCO [11], resulting in poor recognition accuracy for region-specific items including local Egyptian products, Middle Eastern food items, and Arabic-labeled packaging. Currency recognition is similarly limited, failing to accurately identify Egyptian pounds or other regional currencies, particularly when banknotes are worn, folded, or partially visible. From a technical perspective, the application relies on computer vision architectures predating 2020, which lack the inference speed and semantic understanding capabilities of modern state-of-the-art models such as YOLOv8 and vision transformers. This results in higher latency and reduced accuracy compared to contemporary deep learning approaches. Furthermore, while the application offers a hybrid processing model combining on-device and cloud-based inference, the lack of complete offline functionality limits its applicability in connectivity-constrained environments.

\paragraph*{E3rafli by Vodafone (2017)}

E3rafli [10], developed by Vodafone Egypt in 2017, represents one of the few assistive applications specifically designed for the Egyptian market and Arabic-speaking users. The application provides a native Arabic user interface and voice feedback system, addressing a critical linguistic gap in the assistive technology landscape. However, its technical capabilities remain substantially limited compared to global alternatives. The application demonstrates acceptable performance only in two specific domains: color detection and Egyptian currency recognition. These features employ relatively simple computer vision algorithms that can operate effectively on limited training data. In contrast, the application’s general object detection capabilities suffer from low accuracy, likely due to insufficient training data and the use of lightweight, less sophisticated detection models. Scene understanding functionality is similarly restricted, providing only basic environmental descriptions without contextual awareness or safety-relevant information filtering. Document reading and OCR performance are particularly weak, with poor text recognition accuracy for both printed and handwritten Arabic text, limiting the application’s utility for accessing written information. The application is available exclusively on Android devices and employs primarily on-device processing, offering privacy advantages over cloud-dependent alternatives. However, the limited feature set and accuracy constraints render it insufficient as a comprehensive assistive tool for daily independence, necessitating supplementation with additional assistive technologies.

\paragraph*{Comparative Analysis}

Table 2.1 provides a comparative analysis of the key commercial assistive applications, highlighting their strengths and limitations across critical dimensions.

\paragraph*{Table 2.1: Comparison of Commercial Assistive Applications}

\paragraph*{Application}

\paragraph*{Year}

\paragraph*{Platform}

\paragraph*{Arabic Support}

\paragraph*{Processing}

\paragraph*{Key Limitations}

Microsoft Seeing AI iOS only Limited Cloud-based Inaccurate in clutter; cogni- [7] tive overload Google Lookout [8] Android only Limited

\paragraph*{Cloud-based}

Frequent crashes; high battery usage; privacy concerns Envision AI [9] iOS \& Android Weak Hybrid Poor localization; outdated CV models

\paragraph*{E3rafli}

(Voda- Android only fone) [10] Native

\paragraph*{Local}

Limited features; inaccurate object detection

\subsection{Review of Academic Research}

While commercial applications reveal practical limitations in deployment, academic research provides the theoretical foundation and technical innovations necessary to address these challenges. Recent scholarly work has concentrated on three critical dimensions: (1) optimizing state-of-the-art artificial intelligence models for resourceconstrained mobile deployment, (2) developing localization techniques for non-Western languages and cultural contexts, and (3) implementing privacy-preserving architectures that eliminate dependence on cloud infrastructure. This section systematically reviews key contributions in each domain, highlighting their relevance to the development of next-generation assistive systems.

\paragraph*{Image Quality and Interactive Visual Question Answering (VQA)}

A persistent challenge in assistive computer vision is that visually impaired users frequently capture images that are blurry, poorly framed, or partially obscured, resulting from the inherent difficulty of camera alignment without visual feedback. Traditional image captioning systems provide static, single-utterance descriptions that often fail to address the specific informational needs of users in these suboptimal capture scenarios. Huh et al. (2024) conducted a comprehensive user study to address this limitation, proposing a “Long-Form Visual Question Answering” (LF-VQA) framework [1]. Their empirical findings demonstrate that visually impaired users derive significantly greater value from interactive multimodal models that support iterative questioning, allowing them to request clarification, probe specific regions of interest, and obtain contextual information beyond initial descriptions. The study revealed that interactive dialogue-based systems improved task completion rates by 40\% compared to static captioning approaches, particularly for scenarios involving document analysis, navigation decisions, and object identification. This research validates the critical necessity of transitioning from one-way image description to bidirectional conversational assistive agents capable of natural language interaction.

\paragraph*{Arabic Optical Character Recognition (OCR)}

Optical character recognition systems, while achieving accuracy rates exceeding 95\% on printed Latin scripts, exhibit substantially degraded performance when processing Arabic text due to several linguistic and typographic characteristics. Arabic script is inherently cursive, with character shapes varying significantly based on positional context (isolated, initial, medial, or final forms). Additionally, diacritical marks, ligatures, and right-to-left text flow introduce complexities that generic multilingual OCR engines struggle to handle accurately. Ayadi et al. (2025) conducted a systematic investigation of this challenge, proposing specialized deep learning architectures including optimized Generative Adversarial Networks (GANs) specifically trained on Arabic handwritten and printed text datasets [2]. Their experimental results demonstrated that generic multilingual OCR models (such as Tesseract with Arabic language packs) achieve only 68\% character-level accuracy on handwritten Arabic text, whereas their specialized GAN-based architecture achieved 89\% accuracy on the same test set. The study conclusively establishes that reliable Arabic OCR for assistive applications requires dedicated, culturally-aware models trained on region-specific datasets encompassing diverse Arabic fonts, handwriting styles, and document types. This finding has direct implications for developing assistive systems targeting Arabic-speaking populations in the MENA region.

\paragraph*{Localized Currency Recognition}

Financial independence represents a critical component of daily autonomy for visually impaired individuals, requiring reliable identification of local currency denominations. However, global currency recognition models exhibit poor generalization to regional currencies, particularly for nations outside North America and Western Europe. This limitation is exacerbated when banknotes are worn, folded, partially visible, or photographed under variable lighting conditions. Ghanem et al. (2025) addressed this gap by developing a specialized deep learning system for Egyptian currency recognition, constructing a custom dataset of 5,000 images encompassing all current Egyptian pound denominations across varying conditions [3]. Their comparative evaluation demonstrated that transfer learning from global currency datasets achieved only 72\% recognition accuracy on Egyptian notes, whereas their locally-trained convolutional neural network architecture achieved 94\% accuracy under controlled conditions and 87\% accuracy with worn or partially occluded banknotes. The study provides empirical validation for the necessity of region-specific training data and localized model development in assistive applications, particularly for underrepresented currencies and cultural contexts. These findings directly inform the design requirements for currency recognition modules in assistive systems targeting the Egyptian and broader MENA market.

\paragraph*{Real-Time Obstacle Detection (Safety)}

For outdoor navigation assistance, detection latency directly impacts user safety, as delayed obstacle warnings provide insufficient reaction time to prevent collisions. Real-time object detection on mobile devices thus represents a critical technical requirement, necessitating inference speeds exceeding 10 frames per second while maintaining acceptable accuracy. Wang et al. (2024) addressed this challenge by introducing “YOLO-OD,” a modified YOLO (You Only Look Once) architecture specifically optimized for detecting outdoor navigation obstacles including curbstones, poles, stairs, and pedestrian barriers [4]. Their architecture modifications included lightweight feature extractors, optimized anchor configurations for obstacle aspect ratios, and pruned network layers to reduce computational complexity. Experimental evaluation on mobile devices demonstrated that YOLO-OD achieved 15.2 FPS inference speed with 88.3\% mean Average Precision (mAP) for obstacle detection, compared to 8.7 FPS and 91.2\% mAP for standard YOLOv5. The research conclusively demonstrates that single-stage detection architectures offer the optimal trade-off between inference speed and accuracy for real-time mobile deployment, providing sufficient performance for safety-critical navigation assistance while operating entirely on-device without cloud dependency. This work validates the technical feasibility of deploying state-of-the-art object detection models on resource-constrained mobile hardware for latency-sensitive assistive applications.

\paragraph*{Indoor Navigation and Structural Recognition}

Indoor navigation for visually impaired users requires identification of specific structural and architectural elements beyond generic object classes. Critical indoor elements include doorways, staircases, elevators, handrails, chairs, and tables, each requiring precise localization for safe navigation and spatial awareness. Wojtowicz (2025) conducted a comprehensive comparative evaluation of YOLOv8 and YOLOv9 architectures within assistive Mixed Reality (MR) applications, specifically assessing their performance on indoor navigation-relevant object classes [5]. The experimental study employed a custom indoor navigation dataset comprising 12,000 annotated images across diverse indoor environments including offices, homes, and public buildings. Results demonstrated that YOLOv8 achieved 91.7\% precision and 87.3\% recall for detecting doors, 89.2\% precision for staircases, and 85.6\% precision for identifying empty seats. Importantly, YOLOv8 maintained real-time inference speeds of 12.8 FPS on mid-range mobile devices, meeting the responsiveness requirements for assistive navigation applications. The study’s findings validate YOLOv8’s suitability for the “Indoor Mode” of navigation assistance systems, confirming its ability to accurately detect structurally important elements while maintaining computational efficiency compatible with mobile deployment constraints.

\paragraph*{Privacy-Preserving Computer Vision}

The architectural reliance on cloud processing in existing commercial assistive applications creates substantial privacy vulnerabilities, as users must continuously transmit sensitive visual data including images of their home environments, personal documents, financial information, and biometric features to third-party servers. This data exposure contradicts fundamental privacy principles and creates risks of unauthorized access, data breaches, and secondary use of personal information. Martin \& Moon (2025) proposed a comprehensive “Privacy-Preserving Image Captioning” framework that addresses these concerns through a hybrid architecture combining on-device detection with selective encryption [6]. Their approach employs a multi-stage pipeline: (1) lightweight object detection models run locally to identify sensitive visual elements including faces, text, and personally identifiable information; (2) detected sensitive regions are encrypted or redacted using partial image encryption techniques before any data transmission; (3) only privacy-filtered images are transmitted for cloud-based semantic analysis when needed. Experimental evaluation demonstrated that this approach reduces sensitive data exposure by 94\% while maintaining 89\% of the semantic understanding capabilities compared to unrestricted cloud processing. The framework supports implementation of an On-Device Privacy Guard architecture, ensuring that privacy-critical detection and filtering occurs locally before any network transmission, effectively providing user control over data exposure. This research validates the technical feasibility of privacy-preserving assistive systems that balance functional capability with data protection requirements.

\section{Research Gap and Summary}

A systematic analysis of both commercial applications and academic literature reveals a fragmented ecosystem in assistive technology for visually impaired users. While significant progress has been achieved in individual technical domains—including object detection, optical character recognition, and natural language processing— the field lacks integrated, culturally-aware, and privacy-centric systems that address the multidimensional needs of users in non-Western contexts. This section identifies four critical research gaps that current state-of-the-art solutions fail to address comprehensively.

\subsection{The Geo-Linguistic Disparity Gap}

Existing commercial assistive applications and academic research exhibit a pronounced Western-centric development bias, optimizing for English-language contexts and environments typical of North American and Western European settings. Global applications such as Microsoft Seeing AI [7] and Envision AI [9] demonstrate high performance on Latin scripts and Western product recognition but exhibit substantially degraded accuracy when deployed in Middle Eastern contexts. This geo-linguistic disparity manifests across multiple technical dimensions: Linguistic Limitations: Current systems struggle fundamentally with the morphological complexity of Arabic script. The cursive, context-dependent nature of Arabic text, combined with diacritical marks and bidirectional text flow, results in OCR error rates 3–5 times higher than equivalent Latin-script recognition. Furthermore, Arabic dialect variation across the MENA region introduces additional recognition challenges that monolithic language models fail to address. Existing applications provide limited or non-existent support for Arabic voice commands, text-to-speech output with natural Arabic prosody, and culturally appropriate audio feedback mechanisms. Cultural and Environmental Blindness: Object detection models trained on Western-centric datasets such as COCO [11] demonstrate poor generalization to region-specific entities. Critical daily-use items including local food products, Arabic-labeled packaging, regional clothing items, and Egyptian currency denominations remain unrecognized or misclassified. This limitation extends to architectural and environmental elements, as Middle Eastern building structures, street furniture, signage conventions, and urban layouts differ significantly from Western training environments, resulting in degraded navigation assistance accuracy. Dataset Availability: The scarcity of publicly available, annotated datasets for Arabic text in natural scenes, Egyptian products, Middle Eastern architectural elements, and regional currency severely constrains the development of localized models. This data gap perpetuates the Western-centric bias and prevents effective transfer learning for MENA-specific applications.

\subsection{The Obsolescence of Legacy Computer Vision Architectures}

A substantial proportion of deployed assistive applications relies on computer vision models and architectural paradigms developed between 2017 and 2020, predating significant advances in efficient neural network design, transformer-based architectures, and mobile-optimized inference engines. This technological lag creates a critical performance-accuracy-efficiency trade-off that limits practical utility: Latency and Real-Time Performance: Legacy cloud-dependent architectures introduce end-to-end latencies ranging from 800ms to 3 seconds, encompassing image capture, network transmission, cloud inference, and response delivery. For safety-critical navigation scenarios, this latency is unacceptable, as obstacles require detection and warning within 200–300ms to enable appropriate user response. Older on-device models, while offering lower latency, suffer from reduced accuracy due to over-aggressive model compression and outdated architectural designs. Semantic Understanding Limitations: Pre-2023 object detection models typically operate as simple “taggers,” providing class labels and bounding boxes without contextual understanding, spatial relationships, or scene-level semantic reasoning. These systems cannot answer queries such as “Is this intersection safe to cross?” or “Which door leads to the exit?” without explicit rule-based programming. The absence of semantic understanding results in information overload, as systems cannot filter safety-relevant information from inconsequential environmental details. Architectural Inefficiency: Modern architectures including YOLOv8/v10 [15], EfficientDet, and vision transformers offer 2–3× improvements in inference efficiency and 15–20\% accuracy gains compared to legacy models, while requiring comparable or reduced computational resources. The gap in utilizing next-generation 2024–2025 architectures, including Large Multimodal Models (LMMs) [16] for semantic understanding, represents a missed opportunity for substantial capability enhancement in assistive applications.

\subsection{The Privacy-Utility Paradox}

Current assistive technology architectures present users with a forced dichotomy between functional capability and data privacy, requiring them to choose between high-accuracy cloud-based systems that expose sensitive personal information and privacy-preserving offline tools with substantially limited functionality. This paradox manifests across multiple dimensions: Cloud Dependency and Data Exposure: Commercial applications achieving high accuracy rates (Microsoft Seeing AI, Google Lookout, Envision AI) require continuous streaming of visual data to cloud servers for inference. This architecture exposes users to multiple privacy risks: (1) transmission of images containing sensitive personal information, private documents, and home environment details to third-party commercial servers; (2) potential secondary use of collected data for model training, product development, or commercial purposes beyond the stated assistive function; (3) vulnerability to data breaches, unauthorized access, or government surveillance; and (4) lack of user control over data retention, deletion, and usage policies. Limited Offline Capabilities: Existing privacy-preserving approaches that operate entirely offline (such as basic on-device OCR or simple color detection) offer substantially reduced functionality, with accuracy rates 20–40\% lower than cloud-based alternatives. These systems typically support only individual isolated functions rather than integrated multimodal assistance, forcing users to employ multiple applications for comprehensive environmental understanding. Research Gap in Privacy-First Architectures: Academic literature reveals a notable scarcity of research addressing privacy-preserving architectures that maintain competitive accuracy and comprehensive functionality. Specifically, there is a gap in implementing on-device privacy filtering and selective redaction systems that protect sensitive information while enabling cloud-assisted semantic understanding when explicitly authorized by users. Furthermore, insufficient research addresses the optimization of state-of-the-art models for complete on-device execution on mobile hardware, particularly for computationally intensive tasks such as scene understanding and visual question answering.

\subsection{Lack of Interactive Contextual Awareness}

Existing assistive applications predominantly operate as unidirectional communication systems, delivering pregenerated descriptions or simple object labels without supporting natural conversational interaction or contextaware information filtering. Academic research, particularly the work by Huh et al. (2024) [1], establishes that visually impaired users strongly prefer interactive Visual Question Answering (VQA) systems over static image captioning. However, current commercial applications fail to provide this capability due to several limitations: Absence of Dialogue Systems: Most applications provide single-utterance outputs (e.g., “I see a chair, a table, and a person”) without supporting follow-up queries such as “Where is the chair located?” or “Is the path to the door clear?” This limitation prevents users from obtaining specific information relevant to their immediate needs and decision-making contexts. Lack of Intelligent Information Filtering: Existing systems employ indiscriminate description generation, narrating all detected objects regardless of relevance to user safety or current task context. This results in cognitive overload and difficulty extracting actionable information from verbose environmental descriptions. There is a critical need for context-aware filtering that prioritizes safety-critical information (obstacles, hazards, navigation landmarks) over incidental environmental details. Limited Semantic Reasoning: Current applications cannot perform scene-level reasoning or spatial relationship understanding required for complex queries. Questions such as “Can I reach the exit without encountering obstacles?” or “What items are on the table within my reach?” remain beyond the capabilities of existing assistive systems. In conclusion, while individual technical components for superior assistive systems exist in isolation within recent research contributions, no single integrated platform effectively combines localized Arabic AI, real-time safetyoriented processing, privacy-preserving architectures, and interactive contextual awareness into a unified mobile solution. Wareeny is proposed specifically to address this multi-dimensional research gap, providing a comprehensive assistive system tailored to the linguistic, cultural, and privacy requirements of visually impaired users in the Arab world.

\section{Summary of Existing Technologies}

To address the identified research gaps and develop a comprehensive assistive system for visually impaired users in the Arab world, this project strategically leverages a curated suite of state-of-the-art technologies, frameworks, and platforms available as of 2025. This section provides a systematic overview of the key technological enablers that facilitate the development of a responsive, localized, and privacy-aware mobile assistive application. The technology stack is selected based on three primary criteria: (1) technical maturity and proven performance in production environments, (2) compatibility with mobile deployment constraints including computational efficiency and battery consumption, and (3) support for privacy-preserving on-device processing architectures.

\subsection{Cross-Platform Mobile Development Frameworks}

Modern assistive applications face the critical requirement of broad platform accessibility, necessitating support for both iOS and Android ecosystems to maximize user reach across diverse socioeconomic contexts. Flutter [12], developed by Google and released in stable production form in 2018, has emerged as the leading cross-platform mobile development framework for assistive technology applications. Flutter employs the Dart programming language and a reactive framework architecture, rendering user interfaces through a custom graphics engine (Skia) that provides pixel-perfect consistency across platforms without reliance on platform-specific UI components. For assistive technology development, Flutter offers several critical technical advantages. First, it provides native performance characteristics with 60 FPS rendering capabilities, essential for real-time camera preview and visual feedback interfaces. Second, Flutter supports direct integration with platform-specific APIs through platform channels, enabling access to device cameras, sensors, haptic feedback systems, and accessibility services. Third, the framework facilitates seamless embedding of machine learning inference engines including TensorFlow Lite [13] and PyTorch Mobile [14] through Foreign Function Interface (FFI) bindings, enabling on-device AI processing without network dependency. Finally, Flutter’s plugin ecosystem includes pre-built packages for text-to-speech, speech recognition, and camera manipulation, accelerating development of assistive features while maintaining code reusability across platforms.

\subsection{Real-Time Object Detection Architectures}

Real-time object detection represents a foundational capability for assistive navigation and environmental awareness systems. The technical requirements are stringent: detection algorithms must achieve inference speeds exceeding 10 frames per second on mobile hardware while maintaining mean Average Precision (mAP) above 35\% on standard benchmark datasets. The YOLO (You Only Look Once) family of single-stage detection architectures has established dominance in this domain due to their optimal balance of speed and accuracy. YOLOv8 [15], released by Ultralytics in 2023, represents the current state-of-the-art in efficient object detection. The architecture introduces several technical innovations over previous versions: (1) an anchor-free detection head that eliminates the need for predefined anchor box configurations, improving generalization across object scales; (2) a modified CSPDarknet backbone with cross-stage partial connections that reduce computational complexity while maintaining feature extraction capability; (3) a Path Aggregation Network (PANet) feature pyramid that enhances multi-scale object detection; and (4) optimized loss functions combining Classification, Localization, and Distribution Focal Loss for improved training convergence. Empirical evaluation demonstrates that YOLOv8n (nano variant) achieves 37.3\% mAP on the COCO dataset [11] while maintaining inference speeds of 15.2 FPS on mid-range mobile processors with only 6.2 MB model size. This performance profile makes it ideally suited for safety-critical assistive applications requiring detection of dynamic obstacles, traffic signals, indoor structural elements, and in-hand objects with minimal latency. The model’s efficiency enables continuous real-time processing without excessive battery consumption, addressing a critical constraint identified in existing commercial applications such as Google Lookout [8].

\subsection{Large Multimodal Models for Semantic Understanding}

The emergence of Large Multimodal Models (LMMs) in 2023–2024 has fundamentally transformed the landscape of visual understanding and human-computer interaction. Unlike traditional computer vision systems that operate as discriminative classifiers (identifying predefined object categories), LMMs employ generative transformer architectures trained on billions of image-text pairs, enabling semantic reasoning, contextual understanding, and natural language interaction with visual content. OpenAI’s GPT-4 with Vision (GPT-4V) [16], released in September 2023, exemplifies this new paradigm. The model processes images through a vision encoder that generates visual token representations, which are then integrated with text tokens in a unified transformer architecture capable of joint reasoning across modalities. This approach enables several capabilities critical for assistive technology: (1) Visual Question Answering (VQA), allowing users to pose natural language queries about image content and receive contextually relevant responses; (2) scene understanding with spatial reasoning, enabling descriptions of object relationships, spatial layouts, and navigational affordances; (3) document understanding, including extraction and interpretation of text within complex layouts, tables, and forms; and (4) interactive dialogue, supporting multi-turn conversations that allow users to iteratively refine queries and obtain clarification. For assistive applications, LMMs address the fundamental limitation of static image captioning identified in research by Huh et al. (2024) [1], enabling visually impaired users to interrogate their environment through natural conversation. However, current LMM architectures require substantial computational resources (GPT-4V operates exclusively via cloud API), introducing latency and privacy concerns. The strategic deployment of LMMs in the proposed system architecture therefore necessitates a hybrid approach: lightweight on-device models for latencysensitive safety detection, with selective cloud-based LMM invocation for complex semantic understanding tasks when explicitly requested by users and subject to privacy filtering.

\subsection{On-Device Optical Character Recognition Engines}

Optical character recognition represents a critical enabling technology for visually impaired users, providing access to printed documents, product labels, signage, and digital text displays. For assistive applications, OCR engines must satisfy several technical requirements: (1) multilingual support including Arabic script recognition, (2) on-device inference capability to ensure privacy and offline functionality, (3) robustness to varying lighting conditions, text orientations, and image quality, and (4) computational efficiency compatible with mobile deployment. Google ML Kit [17], introduced in 2018, provides an on-device text recognition API leveraging deep learning models trained on diverse text datasets. The framework supports over 100 languages including Arabic, with specialized models for dense text (documents) and sparse text (scene text such as signs and labels). ML Kit employs a two-stage pipeline: text detection using an adapted EAST (Efficient and Accurate Scene Text) detector, followed by text recognition using recurrent neural networks with attention mechanisms. The on-device models require approximately 30 MB storage and achieve inference speeds of 100–300 ms per image on modern mobile devices. For Latin scripts, accuracy exceeds 95\% on printed text, though performance degrades to 70–80\% on handwritten text and complex scene text scenarios. EasyOCR [18], an open-source alternative developed by JaidedAI, provides greater flexibility and customization capability. Built on PyTorch, EasyOCR supports 80+ languages with dedicated models for Arabic script. The framework allows fine-tuning on custom datasets, enabling optimization for region-specific fonts, handwriting styles, and domain-specific terminology. However, EasyOCR’s model sizes (40–60 MB per language pair) and inference times (300–800 ms) present trade-offs against ML Kit’s efficiency. For Arabic OCR specifically, both engines face challenges stemming from Arabic script characteristics discussed in Section 2.2.2. Research by Ayadi et al. (2025) [2] demonstrates that specialized Arabic OCR models trained on regional datasets can achieve 89\% accuracy compared to 68\% for generic multilingual engines. This finding motivates the incorporation of both Google ML Kit for general text recognition and specialized Arabic OCR models for enhanced accuracy on Arabic-language content in the proposed system architecture.

\subsection{Hybrid Cloud-Local Processing Architectures}

Contemporary assistive system design increasingly adopts hybrid architectures that strategically partition computational workloads between on-device (edge) processing and cloud-based inference. This architectural paradigm addresses the fundamental tension between capability, latency, privacy, and resource constraints that characterize mobile assistive applications. In hybrid architectures, latency-sensitive and privacy-critical tasks execute locally on the mobile device using lightweight, optimized models. These tasks typically include real-time obstacle detection, fall detection, traffic light recognition, and personal safety features requiring sub-300ms response times. Conversely, computationally intensive tasks benefiting from large-scale models—such as complex scene understanding, detailed document analysis, and interactive visual question answering—are selectively offloaded to cloud infrastructure when network connectivity is available and user privacy preferences permit. The technical implementation of hybrid architectures requires careful orchestration: (1) intelligent task classification to determine local versus cloud execution based on latency requirements, data sensitivity, and network availability; (2) privacy-preserving preprocessing that filters or encrypts sensitive visual elements before cloud transmission, as proposed by Martin \& Moon (2025) [6]; (3) graceful degradation mechanisms that maintain core functionality when offline or in low-connectivity scenarios; and (4) battery-aware scheduling that balances processing demands against device energy constraints. For the proposed Wareeny system, the hybrid architecture prioritizes user privacy and offline capability while strategically leveraging cloud-based Large Multimodal Models for enhanced semantic understanding when explicitly requested by users and subject to local privacy filtering. This approach addresses the privacy-utility paradox identified in Section 2.3.3, providing comprehensive functionality without compromising data protection principles.


\chaptercover{Chapter 3\\System Analysis and Design}
\chapter{System Analysis and Design}

\section{System Overview and Proposed Concept}

The system overview described in this chapter outlines the architecture and functionality of the AI-powered solution developed to assist visually impaired individuals in navigating both indoor and outdoor environments. The proposed system leverages advanced artificial intelligence techniques, including computer vision, optical character recognition (OCR), object detection, and scene understanding, to analyze visual data captured through a camera-enabled device in real time. The core of the system is built upon a robust AI model trained to recognize obstacles, read text from the surrounding environment, and identify key objects and landmarks that are essential for safe and independent movement. Once the AI model is trained and optimized for performance, it is integrated into a user-friendly application designed specifically for visually impaired users. The application serves as the main interaction platform, providing audio and haptic feedback to convey meaningful information to the user. The AI model processes visual inputs, extracts relevant information such as text, objects, or navigation cues, and delivers clear and timely guidance to help the user understand their surroundings. This integration enhances the user’s ability to move safely and confidently in complex environments, such as streets, buildings, and public spaces. By automating environment perception and information delivery, the system reduces dependence on external assistance and improves overall accessibility. The system model discussed in this chapter focuses on the interaction between the AI components and the application, the flow of data from image capture to user feedback, and how the generated insights are utilized to support real-time decision-making for visually impaired users.

\section{Functional Requirements}

The functional requirements of the Wareeny system are categorized into four main areas of assistance:

\subsection{Indoor Assistance}

FR-IN-1: Indoor Object Detection: The system shall detect indoor objects and estimate distance from them. FR-IN-2: In-Hand Detection: The system shall detect specifically the object that is in the hand of the user, as to not get confused by other objects in the environment. FR-IN-3: Brand Detector: The system shall identify product brands indoors. FR-IN-4: Currency Reader: The system shall detect and read currency denominations. FR-IN-5: Privacy Mode: The system shall detect sensitive data without speaking it aloud. FR-IN-6: Saved Items: The system shall identify user-owned items using stored reference images.

\subsection{Outdoor Assistance}

FR-OUT-1: Outdoor Object Detection: The system shall detect cars, bicycles, and other outdoor obstacles. FR-OUT-2: Outdoor Object Finder: The system shall locate ATMs, benches, trash cans, and other helpful outdoor items. FR-OUT-3: Outdoor Reader: The system shall quickly read shop names, billboards, and street signs. FR-OUT-4: Traffic Light Status: The system shall detect traffic light colors for crossing guidance.

\subsection{OCR Features}

FR-OCR-1: Read Document: The system shall read documents, menus, books, and labels. FR-OCR-2: QA Mode: Instead of just reading a document word for word, the system shall allow users to ask questions about it to get specific responses. FR-OCR-3: Danger Word Alerts: The system shall notify users when dangerous keywords appear, such as “Wet Floor,” “Danger,” or “Caution.”

\subsection{Audio and Haptics Service}

FR-AH-1: Danger Mode: The system shall enable danger mode using voice command, gesture, or phone shake. FR-AH-2: Emergency Auto Call: The system shall automatically call the user’s emergency contact in danger mode. FR-AH-3: Emergency Auto Location: The system shall send the current location of the user to their emergency contact in danger mode. FR-AH-4: Directional Audio: The system shall provide left or right ear cues to indicate object direction.

\section{Non-Functional Requirements}

The non-functional requirements define the quality attributes of the Wareeny system:

\begin{itemize}
  \item Performance: The system shall provide real-time processing with minimal delay.
  \item Accessibility: The system shall offer full control using voice or simple gestures.
  \item Privacy: Local processing is preferred to ensure data protection.
  \item Reliability: The system must function in various indoor and outdoor lighting conditions.
  \item Usability: The system shall provide clear and consistent audio feedback with spatial cues.
  \item Scalability: The system shall support future expansion and improved AI models.
  \item Security: The system shall protect user data and stored personal images.
  \item Power Efficiency: The system shall reduce battery consumption during continuous operation.
\end{itemize}

\section{System Environment}

The proposed system is developed using a combination of modern hardware capabilities and advanced software technologies to ensure efficient, accurate, and real-time assistance for visually impaired users.

\subsection{Hardware Environment}

Table 3.1 summarizes the hardware requirements for the Wareeny system.

\begin{table}[h]
\centering
\caption{Hardware Environment Specifications}
\begin{tabularx}{\textwidth}{|>{\raggedright\arraybackslash}p{0.28\textwidth}|>{\raggedright\arraybackslash}X|}
\hline
\textbf{Component} & \textbf{Specification} \\
\hline
User Mobile Device & A modern smartphone capable of running AI-powered applications. \\
\hline
Processor (CPU) & Multi-core processor (octa-core 2.0 GHz or higher recommended) to support real-time image processing and AI inference. \\
\hline
Memory (RAM) & Minimum 2 GB RAM required; 4 GB or more recommended for smoother execution of deep learning models such as YOLOv8 and PyTorch-based frameworks. \\
\hline
Camera & High-definition rear camera used for capturing images and video streams necessary for object detection, text recognition, and currency identification. \\
\hline
Sensors & Built-in sensors such as the accelerometer and gyroscope to enhance spatial awareness and provide navigation-related feedback. \\
\hline
Connectivity & Active internet connection (Wi-Fi or 4G/5G) required for cloud-based services, including advanced scene understanding and some remote text-to-speech or model-serving workflows. \\
\hline
\end{tabularx}
\end{table}

\subsection{Software Environment}

Table 3.2 presents the software technologies and frameworks utilized in the Wareeny system.

\begin{table}[h]
\centering
\caption{Software Environment Technologies}
\begin{tabularx}{\textwidth}{|>{\raggedright\arraybackslash}p{0.3\textwidth}|>{\raggedright\arraybackslash}X|}
\hline
\textbf{Component/Tool} & \textbf{Purpose and Description} \\
\hline
PyTorch \& OpenCV & Used for model experimentation and computer-vision workflows such as detection, preprocessing, and localized AI prototyping. \\
\hline
YOLOv8 & Employed as the primary real-time object-detection architecture for navigation and stuff-finder scenarios. \\
\hline
EasyOCR \& Google ML Kit & Used for recognizing English and Arabic text from signs, labels, and printed materials in broader research and design considerations. \\
\hline
OpenAI GPT-4 with Vision & Referenced as a cloud-based multimodal reasoning model for advanced scene understanding and complex visual interpretation tasks. \\
\hline
Google Cloud Text-to-Speech (TTS) & Considered in the broader thesis discussion as a high-quality cloud speech option for assistive narration. \\
\hline
Whisper AI & Referenced for potential speech-to-text support in future voice-command and spoken-query interactions. \\
\hline
Flutter (Dart) & Used to develop a cross-platform mobile application compatible with both Android and iOS devices, ensuring accessibility and consistent user experience. \\
\hline
\end{tabularx}
\end{table}

\section{System Architecture}

The system architecture illustrates the high-level design of the Wareeny application, showing how different components interact to provide assistive functionality to the user.

\paragraph*{Figure 3.1: System Architecture Diagram}

\section{Class Diagram}

The class diagram illustrates the structural design of the Wareeny system, depicting the key classes, their attributes, methods, and the relationships between them. The diagram captures the core components of the application, including the user interface controllers, AI processing modules, service classes for object detection, OCR, and audio feedback, as well as data models for storing user preferences and detected objects. The relationships shown—such as associations, compositions, and inheritances—demonstrate how the system components collaborate to deliver seamless assistive functionality. This object-oriented design ensures modularity, maintainability, and scalability, allowing for future enhancements and the integration of additional AI models.

\paragraph*{Figure 3.2: Class Diagram}

\section{Use Case Diagram}

The use case diagram illustrates the interactions between the user and the Wareeny system, showing the various functionalities available to visually impaired users. Figure 3.3 presents the system use cases.

\paragraph*{Figure 3.3: Use Case Diagram}

\section{Sequence Diagrams}

The following sequence diagrams illustrate the flow of interactions for the key features of the Wareeny system.

\subsection{General Flow Diagrams}

\paragraph*{Figure 3.4: General Flow 1: OCR Processing}

\paragraph*{Figure 3.5: General Flow 2: Object Detection Processing}

\paragraph*{Figure 3.6: General Flow 3: Navigation Processing}

\subsection{Indoor Features}

\paragraph*{Figure 3.7: Sequence Diagram: Indoor Object Detection and Distance Estimation}

\paragraph*{Figure 3.8: Sequence Diagram: In-Hand Object Detection}

\paragraph*{Figure 3.9: Sequence Diagram: Brand and Currency Detection}

\paragraph*{Figure 3.10: Sequence Diagram: Saved Items Finder}

\subsection{Outdoor Features}

\paragraph*{Figure 3.11: Sequence Diagram: Outdoor Object Detection}

\paragraph*{Figure 3.12: Sequence Diagram: Outdoor Object Finder}

\paragraph*{Figure 3.13: Sequence Diagram: Outdoor Reader for Signs}

\paragraph*{Figure 3.14: Sequence Diagram: Outdoor Navigation}

\paragraph*{Figure 3.15: Sequence Diagram: Traffic Light Status Detection}

\subsection{OCR Features}

\paragraph*{Figure 3.16: Sequence Diagram: Read Document}

\paragraph*{Figure 3.17: Sequence Diagram: QA Mode (Visual Query)}

\paragraph*{Figure 3.18: Sequence Diagram: Danger Word Alerts}

\subsection{Safety and Emergency Features}

\paragraph*{Figure 3.19: Sequence Diagram: Danger Mode (Hazard Detection)}

\paragraph*{Figure 3.20: Sequence Diagram: Emergency Auto Call and Location Sharing}

\paragraph*{Figure 3.21: Sequence Diagram: Privacy Mode (Secure Shutdown)}

\subsection{Audio and Haptics Features}

\paragraph*{Figure 3.22: Sequence Diagram: Directional Audio}

\section{Entity-Relationship Diagram}

The Entity-Relationship Diagram (ERD) illustrates the database schema and data model of the Wareeny system. The diagram depicts the core entities required to store user information, system preferences, detected objects, saved items, emergency contacts, and usage history. The relationships between these entities—such as one-to-many and many-to-many associations—demonstrate how data flows and persists throughout the application lifecycle. This relational design ensures efficient data storage, quick retrieval of user-specific information, and maintains the integrity of personal data such as saved item references and emergency contact details. The ERD serves as the foundation for implementing a robust and scalable database structure that supports all functional requirements of the assistive system.

\paragraph*{Figure 3.23: Entity-Relationship Diagram}



\chaptercover{Chapter 4\\System Design}
\chapter{System Design}

\section{Introduction}
This chapter presents the design of the current Wareeny application as implemented in the project repository. Wareeny is structured as an accessibility-oriented mobile assistant built with Flutter and centered on camera-driven interaction. In its present form, the system combines a mode-based interface, a vision-language reading pipeline, real-time object-detection flows, spoken feedback, haptic alerts, bilingual localization, and persisted user settings. The implemented design is therefore not a conventional form-based mobile application, but a sensor-driven assistive system in which perception services and accessibility feedback define the core user experience.

\subsection{Design Orientation}
The design follows four main principles. First, interaction must remain usable for blind and low-vision users through two-finger gestures, semantic announcements, spoken prompts, and large touch targets. Second, high-cost perception logic is isolated behind controllers and service abstractions so that interface widgets remain lightweight and reusable. Third, non-realtime understanding tasks are organized around a shared vision-language architecture rather than isolated feature silos. Fourth, real-time navigation and object-finding behavior is delegated to local detection flows in order to reduce latency during active camera use.

\subsection{Implemented Scope}
At the time of this revision, the implemented user-facing flows are the navigation shell, read mode, navigation mode, stuff-finder mode, and settings mode. The mode list also contains a home surface and an in-hand entry, but these do not yet expose dedicated feature controllers comparable to the read, navigation, and stuff-finder modules. Similarly, the internal VLM task model already defines scan, read, in-hand, currency, brand, scene description, and ask-image tasks, but only the reading flow is fully exposed as a finished interaction path in the current application.

\section{Development Tools and Technologies}
\subsection{Flutter and Dart}
Flutter and Dart form the foundation of the application. Flutter provides the cross-platform widget system, theming model, animation support, semantics integration, and localization hooks used across the project. Dart is used for screen composition, stateful feature controllers, dependency registration, bridge abstractions, and AI-service orchestration. This combination is appropriate for Wareeny because it supports fast UI iteration while still allowing native platform integration through method and event channels where lower-level performance is required.

\subsection{Visual Input and Camera Capture}
The application uses two different visual-input paths. For reading, the camera package is used to open the rear camera, display a live preview, and capture a still image on demand. For navigation and stuff-finder mode, the ultralytics\_yolo package provides a streaming camera view tied directly to YOLO inference. This split is architecturally meaningful: reading needs a deliberate still frame for interpretation, whereas navigation and object finding require continuous frame analysis.

\subsection{Vision-Language Reading Stack}
The reading subsystem is built around the VlmAssistant, which accepts an image path and a task definition, then routes the request through a ModelRouter. The router selects between a LocalVlmProvider and a CloudVlmProvider according to task type and runtime availability. Prompt construction is centralized in VlmPrompts so that scan mode and full reading mode can use different instructions while preserving a consistent spoken tone. A SpeechQueue receives streamed text chunks and speaks them sequentially, allowing the application to begin audio output before the full response has completed.

\subsection{Real-Time Object Detection Stack}
Real-time detection functionality is abstracted behind ObjectDetectionSource. The current concrete implementation is UltralyticsYoloDetectionSource, which emits streaming detection frames and selects model assets according to the active feature mode. The repository already contains packaged YOLOv8n model assets for navigation and stuff-finder use on Android and iOS. Detection frames are normalized through DetectionCoordinateMapper before they are interpreted by higher-level logic such as NavigationRiskEngine or StuffFinderController.

\subsection{Localization, Speech, and Preference Persistence}
Flutter localization delegates and a custom AppLocalizations class provide bilingual English and Arabic user-facing strings. Audio output is handled by AudioHapticsBridge, which wraps FlutterTts and configures Arabic-capable voices, speech rate, pitch, and completion callbacks. Haptic feedback is also exposed through the same bridge. Persisted user settings are managed by SettingsController, which stores locale, speech rate, and haptic preference values in SharedPreferences. This design gives the application a lightweight persistence model suitable for personalization without introducing premature database complexity.

\section{System Architecture}
\subsection{Presentation Layer}
The presentation layer consists of NavigationShell and the feature screens it hosts: ReadScreen, NavigationScreen, StuffFinderScreen, and SettingsScreen. Reusable widgets such as CameraTopBar, ModeCarouselBar, TwoFingerSwipeDetector, status chips, overlays, and control bars provide a consistent visual language across modes. The presentation layer is deliberately thin; it captures gestures, displays current state, and delegates feature behavior to dedicated controllers.

\subsection{Application Control Layer}
The application control layer contains stateful controllers that coordinate each feature. VlmReadController manages read-mode state, token streaming, reading status, and playback actions. NavigationController consumes streaming detection frames, applies geometric risk rules, and triggers directional feedback. StuffFinderController manages target selection, detector startup, match counting, and spoken match announcements. SettingsController manages persistent configuration and notifies the UI whenever locale, speech rate, or haptic preferences change.

\subsection{AI Service Layer}
The AI service layer contains the reusable logic that sits underneath the controllers. On the vision-language side, this includes VlmAssistant, ModelRouter, prompt definitions, and the local and cloud provider implementations. On the detection side, it includes the object-detection source abstraction, coordinate normalization, detection models, navigation geometry, and risk scoring. The purpose of this layer is to keep the application extensible: new assistive flows can be added by reusing services instead of duplicating inference logic inside screens.

\subsection{Backend Integration Layer}
The backend integration layer is represented in the current build by CloudVlmProvider. This provider probes a remote \texttt{/health} endpoint, sends multimodal requests to \texttt{/vlm/generate}, and parses server-sent event streams so that tokens can be spoken as they arrive. The backend service itself is external to this repository, but the client-side contract is already encoded in the application. Architecturally, this means the current build is prepared for cloud-assisted reading even when the local runtime is not yet available.

\subsection{Persistence Layer}
The persistence layer is intentionally lightweight in the current build. SharedPreferences is used to store long-lived personalization data, while transient perception data such as streamed text, current navigation assessment, detector frames, and selected temporary targets is held in memory. This means Wareeny currently uses a mixed persistence model: durable key-value settings for personalization and in-memory runtime objects for live camera sessions.

\section{Application Map}
\subsection{Startup and Home Flow}
Application startup begins in bootstrap, where environment configuration is loaded, service registration is performed asynchronously, settings are restored from SharedPreferences, and text-to-speech is initialized. The user is then routed into NavigationShell. From the shell, Home acts as a camera-centered entry surface rather than a form-driven dashboard, and the user can move across modes through the carousel or two-finger swipe gestures.

\subsection{Navigation Mode Flow}
In navigation mode, the shell opens NavigationScreen, instantiates NavigationController, and connects it to a YOLO streaming view. Detection frames are normalized, passed into NavigationRiskEngine, and converted into caution, danger, low-light, or path-clear feedback. Indoor and outdoor submodes are available in the interface, and the same navigation pipeline currently powers both, allowing future mode-specific tuning without redesigning the screen structure.

\subsection{Stuff Finder Mode Flow}
In stuff-finder mode, the user first chooses a target object from a picker that merges custom assistive classes such as wallet, glasses, and keys with COCO object classes. Once a target is selected, StuffFinderController starts the detector, receives streaming frames, filters detections to the selected label, and updates a status chip with the current match state. Spoken feedback announces the selected object and the presence of matches while cooldown logic prevents repetitive output.

\subsection{Read Mode Flow}
In read mode, the user toggles between scan and read submodes. A tap captures a still image from the rear camera. The captured image is converted into a VLM request and routed through the model router. Scan mode is designed for brief context-aware reading, while read mode keeps controls active for a longer spoken response. During reading, streamed text is buffered into speakable sentences and exposed both visually and through text-to-speech.

\subsection{Settings Flow}
Settings mode operates without the live camera and provides language, speech-rate, and haptic-preference controls. The interface uses the same visual style as the rest of the application, but reorganizes it into glass-style configuration cards rather than camera overlays. Every change is persisted immediately so that subsequent launches reproduce the user's preferred accessibility profile.

\section{UI/UX Design}
\subsection{Accessibility Interaction Model}
The interaction model is designed around non-visual control. Two-finger left and right swipes switch top-level modes, while vertical swipes toggle feature-specific submodes in screens such as reading and navigation. Custom semantics actions are attached to several flows so that screen-reader users can invoke mode changes and feature actions without relying solely on visual controls. In this sense, accessibility is not an added layer; it is the primary navigation model of the system.

\subsection{Visual Design Language}
The visual layer uses a dark full-screen camera background with glassmorphism overlays, electric-green highlights, rounded cards, and centrally placed status elements. This design is consistent across the top bar, mode carousel, settings cards, reading controls, and status chips. Even though the target user may not depend on vision, a clean and stable visual layout remains important for caregivers, demonstrations, testing, and hybrid low-vision use.

\subsection{Bilingual Experience}
Wareeny currently supports English and Arabic at the application level. Locale changes propagate through MaterialApp and AppLocalizations, allowing labels, hints, and spoken announcements to adapt to the selected language. The design also benefits from Flutter's bidirectional layout handling, which allows right-to-left presentation behavior when Arabic is selected. This bilingual support is one of the most concrete localization features already implemented in the current build.

\subsection{Feedback Strategy}
Feedback is intentionally multimodal. Visual overlays indicate live status, targets, or reading results; text-to-speech communicates prompts and findings; and haptic signals support confirmation, navigation cues, and danger states. The audio layer also incorporates completion handlers and interruption behavior so that feature controllers can manage reading state more predictably. This combination makes the design more robust than depending on a single output channel.

\section{Database Schema and Persistence Model}
\subsection{Persisted Settings Keys}
\begin{itemize}
  \item \texttt{settings.locale} stores the selected application language as a locale code.
  \item \texttt{settings.speechRate} stores the user's preferred speech output rate as a numeric value.
  \item \texttt{settings.hapticsEnabled} stores whether haptic feedback is enabled or disabled.
\end{itemize}

\subsection{Runtime Data Structures}
The rest of the live application state is intentionally non-persistent. Examples include the currently streamed reading text, active speech queue, latest object-detection frame, latest navigation assessment, selected temporary stuff-finder target, and detector status messages. These values are meaningful only within the current session and therefore remain in controller memory rather than being written to storage.

\subsection{Note on Relational Storage}
The application-based template asks for a relational database schema. In the current Wareeny implementation, however, such a schema would overstate the repository state. The project presently uses key-value preference persistence rather than a relational database engine. Therefore, the academically correct representation for this build is a lightweight persistence schema centered on settings, with relational storage reserved for future phases that may introduce saved items, logs, contacts, or model registries.

\chaptercover{Chapter 5\\System Implementation and Testing}
\chapter{System Implementation and Testing}

\section{Introduction}
This chapter explains how the current Wareeny design has been translated into working application modules. The implementation is mode-oriented and controller-driven. Rather than building all assistive features at once, the repository organizes the work into reusable services that can support multiple assistive tasks over time. The present implementation therefore contains both fully usable user flows and extensibility points for features that are planned but not yet exposed as final screens.

\subsection{Implementation Approach}
The implementation approach is incremental. Shared application concerns such as configuration, service registration, audio output, localization, and persistence are established first. Feature controllers are then added on top of those services so that each mode can own its own logic without duplicating infrastructure. This approach is visible in the separation between shared services such as VlmAssistant and ObjectDetectionSource and feature-specific controllers such as VlmReadController, NavigationController, and StuffFinderController.

\section{System Tools}
\subsection{Frontend and UI Tooling}
Flutter is the main frontend framework, with MaterialApp, animated widgets, semantics, localization delegates, and custom reusable components used throughout the project. Dart is used to implement screens, reusable overlays, controllers, configuration objects, and bridges. The camera package is used in reading mode for still-image acquisition, while standard Flutter state-management patterns such as ChangeNotifier and AnimatedBuilder are used for reactive UI updates.

\subsection{AI and Model Tooling}
The repository currently integrates two AI-oriented tool paths. For reading and future assistive reasoning tasks, a VLM layer is implemented through the local and cloud provider abstractions, routed by ModelRouter and coordinated by VlmAssistant. For real-time navigation and object finding, the ultralytics\_yolo package is used with packaged YOLOv8n model assets. The image package is also used inside the cloud provider to downscale and compress camera captures before remote submission, which reduces transfer and inference overhead.

\subsection{Platform Services and Assets}
Audio output is implemented through FlutterTts inside AudioHapticsBridge, with language-aware configuration and haptic fallbacks. SharedPreferences provides lightweight settings storage. Application model assets are bundled in \texttt{assets/models} and include a general YOLOv8n detector and a separate stuff-finder detector for custom classes. Environment-dependent backend and model identifiers are managed through AppConfig, which distinguishes development and production settings.

\section{Core Features Implemented}
Because Wareeny is an assistive mobile application rather than a CRUD-oriented information system, its implemented core features are better described as functional modes and service flows rather than database transactions. The only true persisted data operations currently implemented are the reading, updating, and reuse of user settings through SharedPreferences. The major implemented functional features are summarized below.

\subsection{Mode-Based Navigation Shell}
NavigationShell is the top-level interaction hub. It speaks a greeting on first launch, exposes the application modes through a reusable carousel, and supports two-finger swipe switching between modes. The shell contains dedicated branches for settings, read mode, navigation mode, and stuff-finder mode, while other mode entries remain structurally visible for future expansion. This gives the application a stable assistive framework even while individual features continue to evolve.

\subsection{Read Mode with Hybrid VLM Routing}
Read mode is implemented through ReadScreen and VlmReadController. The user captures a still image and then chooses either scan mode for brief context-aware reading or read mode for longer spoken reading with controls. The controller creates a VLM request, selects a task type, and sends the request through VlmAssistant. ModelRouter attempts on-device handling first for scan and read tasks, then falls back to cloud processing when the local runtime is unavailable. Streamed token deltas are collected into sentence-sized chunks and spoken through SpeechQueue, while the interface maintains reading controls for stop, repeat, skip, and cancel actions.

\subsection{Real-Time Navigation Mode}
Navigation mode is implemented through NavigationScreen, NavigationController, and NavigationRiskEngine. A YOLO streaming view produces detection frames, which are normalized and passed into a geometric risk model. The controller filters unstable detections through confirmation windows and cooldown rules before generating spoken messages or haptic alerts. It can announce left-side avoidance, right-side avoidance, a frontal stop condition, low-light conditions, or a path-clear state after enough safe frames. Indoor and outdoor navigation submodes are available in the UI, and although both currently share the same detection pipeline, the submode structure prepares the system for later environment-specific tuning.

\subsection{Stuff Finder Mode}
Stuff finder is implemented through StuffFinderScreen and StuffFinderController. The user selects a target from a curated object list that includes both common COCO classes and custom assistive targets such as wallet, glasses, and keys. After selection, the detector is started, the appropriate model path is chosen, and streaming frames are filtered so that only detections matching the selected label remain visible to the feature logic. The controller announces the chosen target, reports matches through speech, and avoids repetitive announcements through a cooldown mechanism.

\subsection{Settings, Localization, and Personalization}
Settings are implemented through SettingsScreen and SettingsController. The user can switch between English and Arabic, adjust speech rate, and enable or disable haptics. These values are persisted immediately and reflected across the application through animated rebuilds. This feature is important because Wareeny is not only an inference tool; it is also a personalized assistive interface whose output behavior must adapt to the user's preferred language and feedback style.

\subsection{Audio and Haptic Delivery}
Audio and haptic behavior is handled centrally so that multiple features can share a common feedback layer. AudioHapticsBridge initializes speech early, selects an Arabic-capable voice when available, applies speech-rate updates, and supports interruption or completion callbacks. Haptics are exposed through semantic categories such as danger, obstacle, navigation cue, and confirmation. This shared delivery layer simplifies controller design because feature modules do not need to manage raw TTS configuration independently.

\section{User Interface Organized and Documented}
\subsection{Reusable Widgets}
The interface is organized around reusable widgets rather than large monolithic screens. Examples include CameraTopBar, ModeCarouselBar, CameraStatusView, CameraSideActions, OcrControlsBar, and multiple overlay widgets. This organization improves maintainability and supports consistent styling across modes. It also makes documentation easier, because each widget represents a clearly bounded responsibility within the overall interface.

\subsection{Screen Composition and Overlays}
Each major feature screen is built as a stack over a full-screen camera layer or a settings surface. Overlays are then placed for branding, status, target selection, reading state, debug information, or mode switching. This layered composition is particularly suitable for assistive camera applications because it preserves a stable visual hierarchy while allowing live perception information to appear without replacing the underlying preview.

\subsection{Accessibility Semantics}
Accessibility semantics are explicitly integrated into several parts of the interface. Mode switching exposes custom semantics actions, important status values are marked as live regions, target selectors and controls are labeled for screen-reader use, and announcements are spoken whenever major mode changes occur. This means the UI documentation for Wareeny is not purely visual; the semantics model is part of the user-interface implementation itself.

\section{Source Code Organized and Documented}
\subsection{Bootstrap and Dependency Injection}
Application startup is centralized in bootstrap, which initializes widget binding, loads configuration, restores settings, builds the service registry, and eagerly prepares audio services. ServiceRegistry then provides a single place for wiring the settings controller, bridge classes, VLM assistant, and object-detection source. This is a strong organizational choice because it keeps cross-cutting dependencies out of feature screens.

\subsection{Feature Controllers}
Feature-specific behavior is encoded in controllers rather than view code. VlmReadController manages the reading state machine. NavigationController manages streaming detection analysis and feedback timing. StuffFinderController manages target-based detection and spoken announcements. SettingsController manages persistent preferences. This controller-based structure improves readability, testability, and future maintainability because business logic is separated from widget rendering.

\subsection{Shared AI and Detection Services}
Shared services such as VlmAssistant, ModelRouter, LocalVlmProvider, CloudVlmProvider, ObjectDetectionSource, and NavigationRiskEngine serve as reusable infrastructure. They allow multiple modes to share the same inference foundations without coupling controllers directly to backend details or model implementations. This is especially important in assistive software, where new features often extend existing perception capabilities rather than starting from scratch.

\subsection{Current Implementation Boundaries}
Although the architecture is broader than the current exposed UI, some boundaries remain important for accurate documentation. The local VLM provider depends on a native on-device runtime that is not yet fully wired on all platforms. The cloud VLM provider expects a backend service external to this repository. Ask-image, currency, brand, and dedicated in-hand flows are represented at the task or mode level but are not yet delivered as finished screens. The current persistence model is still key-value based rather than relational. These boundaries should be described explicitly so that the thesis remains precise and credible.

\section{Test Plan and Test Cases Prepared}
\subsection{Functional Test Cases}
\begin{itemize}
  \item TC-01: Verify that application startup restores persisted language, speech-rate, and haptic preferences and presents the navigation shell correctly.
  \item TC-02: Verify two-finger gesture mode switching and correct semantic announcement of the newly active mode.
  \item TC-03: Verify read-mode scan behavior from still-image capture through VLM response and spoken completion.
  \item TC-04: Verify full read behavior, including token streaming, stop, repeat, skip, and cancel controls.
  \item TC-05: Verify navigation-mode warning generation for left, right, frontal danger, low-light, and path-clear conditions.
  \item TC-06: Verify stuff-finder target selection, detector startup, match counting, and cooldown-based speech behavior.
\end{itemize}

\subsection{Integration and Fallback Cases}
\begin{itemize}
  \item TC-07: Verify local-provider unavailability and correct cloud fallback for reading tasks.
  \item TC-08: Verify unavailable-state messaging when neither local nor cloud reading can serve the request.
  \item TC-09: Verify model-path switching between the general YOLOv8n asset and the custom stuff-finder asset.
  \item TC-10: Verify settings persistence across application restarts.
\end{itemize}

\subsection{Current Automated Coverage}
The repository includes widget-test scaffolding for shell-level behavior, but the broader assistive flows remain primarily dependent on functional and device-based verification because they rely on camera input, platform views, speech, haptics, model assets, and backend availability. Consequently, the current test plan is stronger as an integration-testing specification than as a fully automated test suite.

\section{Testing the Core Features with Sample Data}
\subsection{Reading Samples}
The reading flow can be functionally evaluated with printed signs, labels, short notices, and multi-line documents captured through the rear camera. The most relevant qualitative checks are whether scan mode produces a short context-aware answer, whether read mode maintains a longer spoken output, and whether the application continues to offer playback controls until the speech queue has drained.

\subsection{Navigation Samples}
Navigation mode is most appropriately tested with live scenes rather than static benchmark files because its logic depends on streaming frames, confirmation windows, and repeated observations. Useful sample scenes include a clear corridor, a centrally placed obstacle, a left-offset obstacle, a right-offset obstacle, and low-light environments. These scenarios verify that geometric risk rules and cooldown timing behave as intended.

\subsection{Stuff-Finder Samples}
Stuff finder can be tested with both custom targets and common COCO targets. For example, keys, glasses, wallet, bottle, chair, and cell phone form a practical sample set because they exercise both the custom-target path and the general object list. The main functional checks are correct target selection, correct match filtering, appropriate status updates, and non-repetitive audio reporting.

\subsection{Settings Persistence Samples}
Settings mode can be tested by changing the language, speech rate, and haptic toggle, then restarting the application to confirm that values are restored correctly. This is the most direct sample-data validation for the current persistence layer because the stored data model is intentionally limited to user preferences.

\subsection{Current Testing Status}
At the current repository stage, the strongest claims that can be made are functional and architectural rather than benchmark-heavy. Wareeny already supports meaningful mode-level validation for reading, navigation, stuff finding, and settings persistence. However, comprehensive latency reporting, formal accuracy studies, multi-device benchmarking, and large-scale user evaluation remain future work and should not be presented as completed outcomes in this thesis version.

\chaptercover{Chapter 6\\Conclusions and Future Work}
\chapter{Conclusions and Future Work}

\section{Introduction}
This chapter summarizes the current state of Wareeny as represented by the implemented repository and identifies the most appropriate future directions. The purpose of this conclusion is not to describe an imagined final product, but to evaluate the present application honestly as an assistive mobile system that has already achieved meaningful architectural and functional progress.

\section{Project Summary}
\subsection{Functional Summary}
Wareeny currently provides a mode-based assistive shell with working navigation, stuff-finder, read, and settings flows. The reading module uses a hybrid vision-language architecture to convert captured images into spoken responses through scan and read submodes. The navigation module uses streamed object detection, path geometry, and risk-based rules to produce caution and danger feedback. The stuff-finder module filters streamed detections to a user-selected target and reports matches. The settings module persists language, speech-rate, and haptic preferences across sessions.

\subsection{Architectural Summary}
From an architectural perspective, the project has progressed beyond a basic prototype. The codebase already separates presentation, feature control, shared AI services, detection abstractions, bridge-based platform services, and persisted settings. It also encodes a clear distinction between realtime local vision flows and higher-level VLM-based understanding tasks. This separation is important because it gives the project a credible path for expansion without requiring a complete redesign.

\section{Project Contributions}
\subsection{Accessibility-Centered Interaction Model}
One important contribution of the project is its explicit accessibility orientation at the interaction level. Two-finger gestures, semantic announcements, spoken greetings, live-region updates, language-aware speech, and configurable haptic behavior are not auxiliary additions; they are part of the core interaction contract. This makes Wareeny more appropriate for assistive use than a system that simply adds speech after completing a primarily visual interface.

\subsection{Hybrid Reading Architecture}
A second contribution is the introduction of a shared reading architecture built around VlmAssistant, routed providers, prompt specialization, and streamed speech delivery. This design is more extensible than a single hardcoded OCR pipeline because it supports multiple task types and explicit local-versus-cloud routing. Even where some provider paths still depend on external runtime availability, the architectural contract itself is already defined and implemented.

\subsection{Real-Time Local Detection Integration}
A third contribution is the integration of real-time YOLO-based detection into mode-specific controllers for navigation and object finding. Rather than merely displaying raw detections, the application converts them into assistive feedback through geometry, target filtering, cooldown logic, and screen-reader-compatible UI layers. This makes the object-detection component practically useful within the larger assistive application context.

\subsection{Extensible Service Structure}
A fourth contribution is the service-oriented organization of the codebase. The current repository already includes clear extension points for backend-assisted VLM use, on-device local runtime integration, battery and thermal bridges, file-download workflows, and additional perception tasks such as brand, currency, and ask-image reasoning. As a result, the project can grow feature-by-feature while preserving a coherent architecture.

\section{Conclusion and Future Work}
\subsection{Conclusion}
In conclusion, Wareeny currently stands as a serious assistive-application foundation rather than a superficial proof of concept. The system already demonstrates a reusable accessibility shell, a hybrid reading pipeline, real-time local detection flows, persisted personalization, and bilingual interaction support. At the same time, the academically correct interpretation of the repository is that some major extensions remain under development and should be reported as future work rather than completed functionality.

\subsection{Near-Term Future Work}
\begin{itemize}
  \item Expose dedicated user flows for the VLM tasks that already exist architecturally, especially in-hand, ask-image, brand, currency, and scene-description interactions.
  \item Complete production-grade integration of the local on-device VLM runtime so that reading and related tasks can operate privately and offline when the model is installed.
  \item Stabilize and validate the remote backend contract used by the cloud VLM provider, including deployment, reliability, and secure request handling.
  \item Expand automated and device-based testing so that navigation, stuff-finder, and reading behaviors can be evaluated across multiple phones and environments.
  \item Strengthen persistence beyond SharedPreferences when the project begins storing richer entities such as saved items, interaction history, or model metadata.
\end{itemize}

\subsection{Long-Term Future Work}
Longer-term work may include richer emergency and safety workflows, stronger Arabic-domain evaluation, more diverse object-finder models, thermal and battery-aware adaptation, and formal user studies with blind and low-vision participants. These extensions would move Wareeny from its current strong foundational state toward a broader and clinically discussable assistive platform. For this thesis version, however, the most accurate and professional conclusion is that the project has successfully implemented a credible architectural and functional base on which these future additions can be built.


\chaptercover{Bibliography}
\addcontentsline{toc}{chapter}{Bibliography}
\begin{thebibliography}{99}
  \bibitem{ref1} M. Huh et al., ``Long-Form Answers to Visual Questions from Blind and Low Vision People,'' arXiv preprint arXiv:2408.06303, 2024. [Online]. Available: \url{https://arxiv.org/abs/2408.06303}
  \bibitem{ref2} M. Ayadi et al., ``Empowering Accessibility in Handwritten Arabic Text Recognition for Visually Impaired Individuals through Optimized Generative Adversarial Network (GAN) model,'' \textit{Journal of Disability Research}, vol. 4, no. 1, 2025. [Online]. Available: \url{https://www.scienceopen.com/hosted-document?doi=10.57197/JDR-2024-0110}
  \bibitem{ref3} A. M. Ghanem, H. A. Youness, M. Wahba, and H. M. Abdelaal, ``Recognizing Egyptian currency for people with visual impairment using deep learning models,'' \textit{Scientific Reports}, vol. 15, no. 1, 2025. [Online]. Available: \url{https://www.researchgate.net/publication/396095205}
  \bibitem{ref4} W. Wang et al., ``YOLO-OD: Obstacle Detection for Visually Impaired Navigation Assistance,'' \textit{Sensors}, vol. 24, no. 23, p. 7621, 2024. [Online]. Available: \url{https://www.mdpi.com/1424-8220/24/23/7621}
  \bibitem{ref5} K. Wojtowicz, ``The Impact of YOLOv8 and YOLOv9 Neural Network Architectures on the Quality and Speed of Object Detection in Assistive VR/MR Applications for Visually Impaired Users,'' \textit{European Research Studies Journal}, vol. 28, no. 4, pp. 1044--1057, 2025. [Online]. Available: \url{https://ersj.eu/journal/4157}
  \bibitem{ref6} A. D. Martin and I. Moon, ``Privacy-Preserving Image Captioning with Partial Encryption and Deep Learning,'' \textit{Mathematics}, vol. 13, no. 4, p. 554, 2025. [Online]. Available: \url{https://www.mdpi.com/2227-7390/13/4/554}
  \bibitem{ref7} Microsoft Corporation, ``Seeing AI,'' 2017. [Online]. Available: \url{https://www.microsoft.com/en-us/ai/seeing-ai}
  \bibitem{ref8} Google LLC, ``Lookout by Google,'' 2019. [Online]. Available: \url{https://blog.google/products/android/lookout/}
  \bibitem{ref9} Envision B.V., ``Envision AI,'' 2018. [Online]. Available: \url{https://www.letsenvision.com/}
  \bibitem{ref10} Vodafone Egypt, ``E3rafli: Mobile Application for the Visually Impaired,'' 2017. [Online]. Available: \url{https://www.vodafone.com.eg/en/e3rafli}
  \bibitem{ref11} T.-Y. Lin, M. Maire, S. Belongie, J. Hays, P. Perona, D. Ramanan, P. Dollar, and C. L. Zitnick, ``Microsoft COCO: Common Objects in Context,'' in \textit{European Conference on Computer Vision (ECCV)}, 2014, pp. 740--755. [Online]. Available: \url{https://arxiv.org/abs/1405.0312}
  \bibitem{ref12} Google LLC, ``Flutter - Build apps for any screen,'' 2024. [Online]. Available: \url{https://flutter.dev/}
  \bibitem{ref13} Google LLC, ``TensorFlow Lite - ML for Mobile and Edge Devices,'' 2024. [Online]. Available: \url{https://www.tensorflow.org/lite}
  \bibitem{ref14} Meta Platforms, Inc., ``PyTorch Mobile,'' 2024. [Online]. Available: \url{https://pytorch.org/mobile/}
  \bibitem{ref15} G. Jocher, A. Chaurasia, and J. Qiu, ``Ultralytics YOLOv8,'' 2023. [Online]. Available: \url{https://github.com/ultralytics/ultralytics}
  \bibitem{ref16} OpenAI, ``GPT-4 with Vision System Card,'' 2023. [Online]. Available: \url{https://openai.com/research/gpt-4v-system-card}
  \bibitem{ref17} Google LLC, ``ML Kit - Machine Learning for Mobile Developers,'' 2024. [Online]. Available: \url{https://developers.google.com/ml-kit}
  \bibitem{ref18} JaidedAI, ``EasyOCR - Ready-to-use OCR with 80+ languages,'' 2024. [Online]. Available: \url{https://github.com/JaidedAI/EasyOCR}
  \bibitem{ref19} World Health Organization, ``World Report on Vision,'' Geneva, Switzerland, 2019. [Online]. Available: \url{https://www.who.int/publications/i/item/9789241516570}
\end{thebibliography}
\end{document}
